\section{Case study: stroke triage across three hospitals}
\label{sec:case-study}

The healthcare demo in the Dagents repository runs a federated study on real,
de-identified patient records. It shows the principles of \cref{sec:design} at
work and measures what the Guard costs in accuracy. The demo's models are not
clinically validated, and the demo is not a medical device.

\subsection{Data and setting}

The records come from the International Stroke Trial (IST-1), which enrolled
19,435 patients with acute ischaemic stroke between 1991 and
1996~\citep{ist1997lancet,sandercock2011ist}. The trial's database is published
under the Open Data Commons Attribution License~\citep{sandercock2011istdata}.
The demo uses 1,014 patients from three hospitals: one in the United Kingdom
(524 patients), one in Italy (292) and one in Switzerland (198). The task is to
flag, at arrival, patients at high risk of death or dependency six months
later. The hospitals agree on a feature contract of 13 fields, such as age band,
level of consciousness, limb weakness, hours since onset and systolic blood
pressure, and the six-month outcome is the label.

Each hospital has its own store and its own LMA, and the study has one GMA. The
application contributes one extension: the feature contract, the data
classification, the condition pack with its approved purposes and release
gates, and two model kinds (the current points rule and a logistic regression).
The classification marks four fields as high sensitivity (hours since onset,
systolic blood pressure, atrial fibrillation and the outcome), nine as medium
(such as age band and sex) and one as low (whether symptoms were noticed on
waking). It also marks the model update as high sensitivity and sets the
minimum cohort to 20.

\subsection{Governance decisions}

\Cref{tab:probe} shows five decisions of one hospital's Guard on the same three
fields. The verified requester scores 0.97 on the KYU scale (high trust); the
unverified one scores 0.445 (moderate trust).

\begin{table}[h]
\centering
\small
\begin{tabularx}{\linewidth}{@{}Y l l l l@{}}
\toprule
Request & \code{systolic\_bp} (high) & \code{age\_band} (medium) & \code{onset\_on\_waking} (low) & Decision \\
\midrule
Verified, row & generalize:1 & allow\_full & allow\_full & narrow \\
Unverified, row & redact & generalize:1 & allow\_full & narrow \\
Verified, table, 25 subjects & aggregate\_only:20 & aggregate\_only:20 & allow\_full & narrow \\
Verified, table, 5 subjects & refuse & refuse & refuse & deny \\
Verified, model update & clip\_contribution:1 & clip\_contribution:1 & clip\_contribution:1 & narrow \\
\bottomrule
\end{tabularx}
\caption{Decisions of the Guard at one hospital, from the demo's governance
endpoint. The table request for 5 subjects is denied because 5 is below the
minimum cohort of 20.}
\label{tab:probe}
\end{table}

\subsection{The study}

The demo's study runs eight rounds in order, each sent by the GMA to every
hospital's LMA: one analytics round, one round that evaluates the model in use,
five training rounds, and one round in which each site validates the candidate
on its own records. The five training rounds lower the candidate's loss from
0.693 to 0.496. The GMA then evaluates five release gates:
\begin{itemize}
  \item discrimination: AUC 0.786, required at least 0.7 (pass);
  \item improvement on the model in use: AUC 0.786, required at least the
        current 0.777 plus a margin of 0.02 (fail, blocking);
  \item subgroup fairness: AUC gap 0.022, required at most 0.1 (pass);
  \item sensitivity at the alert threshold: 0.922, required at least 0.8
        (pass);
  \item alert burden: 853 alerts per 1,000 patients, required at most 750
        (fail).
\end{itemize}
The decision is to reject the candidate, and the hospitals keep the model in
use. The candidate is better than the model in use, but not by the margin the
consortium set. The run has no randomness, so it repeats exactly on the
committed data. The candidate is trained and measured on the same three
hospitals, which raises its AUC by about 0.02 compared with unseen data.

\subsection{What crosses, and what stays}

\begin{table}[h]
\centering
\small
\begin{tabularx}{\linewidth}{@{}Y Y@{}}
\toprule
Crosses to the coordinator & Stays at the hospital \\
\midrule
Record counts, and an outcome rate when the planner permits a table request
for the outcome at that many labelled records (at least 20) & The records,
including every outcome label \\
Each round's aggregate metrics, such as AUC, when \code{before\_send} permits
the round's contribution & The guarded records the round's code trains on \\
The norm and digest of the clipped model update (the update itself goes to a
write-only store) & The audit chain of the hospital's Guard decisions \\
Whether the audit chain verifies, and how many decisions it holds & The
decisions in that chain \\
\bottomrule
\end{tabularx}
\caption{What crosses from a hospital to the coordinator in the demo.}
\label{tab:crosses}
\end{table}

\subsection{The cost of the Guard}
\label{sec:cost}

Two high-sensitivity fields, hours since onset and systolic blood pressure, are
generalised to the nearest 10 when a site reads its records for a round, so the
round's code never sees the exact values (P8). We measured the cost in two ways
on the committed data, with the planner built from the repository:
\begin{itemize}
  \item \textbf{Training.} We ran the pilot twice: once with the committed
        classification, and once with the two fields lowered to medium
        sensitivity, so that the Guard returns them unchanged to the trusted
        runner. The candidate trained on exact values reaches a size-weighted
        AUC of 0.7864; the candidate trained on guarded values reaches 0.7859.
        The cost is 0.0005 AUC.
  \item \textbf{Scoring.} The candidate trained on guarded values scores 0.7594,
        0.8117 and 0.8185 at the three hospitals on their exact records, and
        0.7593, 0.8115 and 0.8184 on the guarded records it was validated on.
\end{itemize}
To change this cost, a consortium changes the classification, as in the first
run above; no code path changes (P7). These numbers measure one model on one
dataset and do not show that the cost is small in general.

\subsection{Other consumers}

A second demo, which turns questions in plain English into SQL, uses the same
planners for source validation, pipeline planning and workload compilation, but
no federation. A fourth hospital, which is not in the study at the start, can
join a running study on one machine: the core service renders a Compose file
for its LMA and the API starts it. These show that the planner and Guard are
not specific to one application, but they are demos run by the authors, not
deployments by independent users.
